\begin{table}[ht!]
\centering
\caption{What the exposure score is built from: the 2019 occupation codes of the employer's own incumbents}
\label{tab:occ_coverage}
\scriptsize
\setlength{\tabcolsep}{5pt}
\begin{tabular}{lrrrr}
\toprule
 & Incumbents & Coded from & Coded, with & Scored \\
 & & 2019 alone (\%) & carry-back (\%) & (\%) \\
\midrule
31--34 & 474,587 & 91.6 & 94.7 & 94.3 \\
35--40 & 672,986 & 92.6 & 95.3 & 95.1 \\
41--49 & 1,063,594 & 93.7 & 96.0 & 95.8 \\
50 and over & 1,907,832 & 90.9 & 93.1 & 92.9 \\
\midrule
All incumbents, 31--69 & 4,118,999 & 92.0 & 94.4 & 94.2 \\
\addlinespace[3pt]
\midrule
\multicolumn{5}{l}{\textit{Employers, and the floor}} \\
With an incumbent aged 31 to 69 & 294,887 & & & \\
\quad scored at a floor of one person-month & 271,047 & & & \\
\quad scored at a floor of three person-months & 267,605 & & & \\
\quad scored at a floor of five person-months (reported) & 262,089 & & & \\
\addlinespace[3pt]
\midrule
\multicolumn{5}{l}{\textit{What the scoring level changes}} \\
\multicolumn{4}{l}{Employers moving quartile against four-digits-where-usable (\%)} & 9.7 \\
\multicolumn{4}{l}{\quad rank correlation of the two firm scores} & +0.977 \\
\multicolumn{4}{l}{Employers moving quartile against four-digits-only (\%)} & 9.8 \\
\multicolumn{4}{l}{\quad rank correlation of the two firm scores} & +0.976 \\
\bottomrule
\end{tabular}
\begin{minipage}{0.92\textwidth}\footnotesize\vspace{4pt}
Incumbents are the employees aged 31 to 69 on the employer's 2019 payroll, the workers the score averages over; no younger worker is scored and no code dated after the 2019 register file enters. A code is taken from the 2019 register where it holds one and from 2018, 2017, 2016 or 2015 otherwise, which is the carry-back column; 2.5 per cent of the resolved codes come from one of those earlier years. Scored is the share the DAIOE index can price after the merge. The reported arm scores every resolved incumbent from the three-digit book: in this delivery every 2019 code carries four digits, so the coarsening is a choice and not a gap, and the last block reports what the alternatives change. A forward arm that also reads 2020 and 2021 resolves 96.3 per cent of incumbents; it is reported in the replication export and is never the score.
\end{minipage}
\end{table}
